\section{Driving Fluent from Claude Code on the DLL remote PCs}
\label{app:claude}

The study-A runs were driven by Claude Code through an MCP server connected to the running Fluent
session. This is optional: every step of this report can be done in the Fluent GUI. The set-up is in
\file{CFD/mcp/} and follows 2026\_C\_FD\_REPORT-2\_MCP-FLUENT, with two differences: the MCP server is
\file{CFD/mcp/server.py} of this repository (the guide does not give the URL of its own), and the
Claude Code home is moved to the network drive with the \texttt{CLAUDE\_CONFIG\_DIR} variable, because
the guide's junction cannot point to a network drive.

\paragraph{Set-up, every session.}
\begin{enumerate}
\item Put the repository in the Documents folder of the network home, which survives the reset of
\texttt{C:\textbackslash} (clone it from there, or download the branch as a ZIP from GitHub).
\item In a normal PowerShell window run
\begin{center}
\small\texttt{powershell -ExecutionPolicy Bypass -File <repo>\textbackslash CFD\textbackslash mcp\textbackslash bootstrap-vm.ps1}
\end{center}
Each step is skipped when its result is already there: the first run takes about 10 minutes, later runs
seconds. It finds the newest Fluent installation and Python, installs PyFluent 0.43 and \texttt{mcp} 2.x
from the wheel cache on the network drive, installs or restores Claude Code, writes \file{.mcp.json}
(without a BOM) and starts Claude Code in the repository.
\item Start Fluent yourself (Workbench or the Fluent launcher, working directory \file{CFD/runs}). In
its console type \texttt{define/beta-feature-access} (answer yes), then \texttt{server/start-server}.
Fluent 2024~R2 refuses \texttt{start-server} until the beta features are on. Fluent writes a
\file{server_info*.txt} file.
\item Ask Claude Code to connect to that file (\texttt{connect\_to\_fluent}). The \texttt{launch\_fluent}
tool can also start a new session, but it always picks the newest installed version (2026~R1 on these
PCs), not 2024~R2.
\end{enumerate}

\paragraph{Restarting.}
\begin{itemize}
\item \textbf{Claude Code only:} type \texttt{/exit}, then run the setup script again, or start it by hand
with \texttt{claude --continue} from the repository folder. \texttt{--continue} reopens the last
conversation; inside Claude Code, \texttt{/resume} picks an older one. The history is kept on the network
drive, so it survives the restart. Fluent can stay open.
\item \textbf{MCP connection only:} type \texttt{/mcp}, choose \texttt{fluent} and \gui{Reconnect}.
Then ask Claude to run \texttt{connect\_to\_fluent} again, with the \file{CFD/runs} folder. PyFluent takes up to a
minute to load.
\item \textbf{Fluent restarted:} type \texttt{server/start-server} again in the new Fluent console, then
reconnect. Each Fluent session writes its own \file{server_info} file, and Claude uses the newest one in
the folder.
\end{itemize}

\paragraph{Tools.} The server offers 13 tools. Session: \texttt{launch\_fluent},
\texttt{connect\_to\_fluent}, \texttt{exit\_fluent}. Files: \texttt{read\_case},
\texttt{read\_case\_and\_data}, \texttt{write\_data}. Solver: \texttt{initialize}, \texttt{iterate},
\texttt{execute\_tui}, \texttt{scheme\_eval}. Reports: \texttt{get\_case\_summary},
\texttt{get\_residuals\_summary}, \texttt{screenshot}. The tool time-out is set to one hour so that
long \texttt{iterate} calls finish.

\paragraph{Things to know.}
\begin{itemize}
\item The server drives the solver only; geometry and meshing stay manual.
\item Console commands with keyword menus are easy to get wrong. Let Claude Code run iterations and read
results, but enter boundary values in the panels and check them with \texttt{/file/write-bc}.
\item One \texttt{iterate} call hung with Fluent waiting at its ``number of iterations'' prompt;
\texttt{/solve/iterate n} through \texttt{execute\_tui} works.
\item The copy of Claude Code on the PC is a separate installation: it starts without the context of
the sessions on a laptop, so give it the case description (this report) at the start.
\end{itemize}
